\chapter{Conclusão}\label{cap:conclusao}

Em diferentes contextos e para variados públicos-alvo, um \textit{software} interativo e de fácil usabilidade pode assumir uma importância fundamental no processo de aplicação da metodologia ABA. Ele precisa ser acessível tanto para pessoas sem conhecimento técnico avançado, quanto para aquelas não familiarizadas com softwares complexos. A fim de permitir que qualquer indivíduo possa utilizá-lo facilmente. 

Essa ferramenta visa oferecer uma gestão mais eficaz de dados e, além disso, apresentar uma abordagem alternativa para os treinos por tentativas discretas. Mediante isso, adaptando-se às necessidades da criança, independentemente do nível do TEA. Tão logo, o software também se ajusta às preferências de trabalho dos terapeutas, oferecendo diversas formas de realizar os treinos.

Acredita-se que a utilização do software possa potencialmente tornar as sessões de terapia mais eficientes, possibilitando uma comunicação mais facilitada e potencialmente fornecendo um canal mais ágil para \textit{feedback} com a família do paciente.

\section{Trabalhos Futuros}

Para futuras pesquisas, há expectativas de aprimoramento nos treinos do programa expressivo, visando uma funcionalidade que possibilite a captação e validação da voz da criança como um critério adicional de sucesso. Além disso, há previsão de implementação de novas funcionalidades na aplicação, ampliando ainda mais suas capacidades e utilidades. Espera-se também realizar estudos para adaptar o software a crianças de diferentes faixas etárias, considerando variações no desenvolvimento e nas necessidades de aprendizagem.

